\documentclass[11pt, letterpaper]{article}

% Essential packages
\usepackage[margin=1in]{geometry}
\usepackage{booktabs} % For professional tables
\usepackage{graphicx} % For inserting lab photos or plots
\usepackage{enumitem} % For customizable lists
\usepackage{hyperref}
\usepackage{fancyhdr}
\usepackage{tcolorbox} % For highlighting issues/troubleshooting

% Header and Footer configuration
\pagestyle{fancy}
\fancyhf{}
\fancyhead[L]{Experimental Lab Notes}
\fancyhead[R]{\today}
\fancyfoot[C]{\thepage}

% Document Information
\title{\textbf{Lab Notebook: [Experiment Name]}\\ \large Subtitle or Project Phase}
\author{Tyler Celotta \\ Reed College}
\date{\today}

\begin{document}

\maketitle
\thispagestyle{fancy} % Applies header/footer to the first page

\section{Objective}
% Briefly state the goal of today's experiment.
To prototype and evaluate the pneumatic response of a silicone-based soft actuator for use in modular prosthetic grip assistance.

\section{Materials \& Equipment}
\begin{itemize}[noitemsep]
    \item Ecoflex 00-30 Silicone (Part A and B)
    \item 3D-printed PLA molds
    \item Syringe and pneumatic tubing (4mm OD)
    \item Digital pressure gauge (range: 0-100 psi)
    \item [Add other materials here]
\end{itemize}

\section{Procedure}
\begin{enumerate}
    \item \textbf{Preparation:} Mixed 50g of Part A and 50g of Part B. Degassed in a vacuum chamber for 10 minutes.
    \item \textbf{Curing:} Poured mixture into the 3D-printed mold and allowed to cure at room temperature for 4 hours.
    \item \textbf{Testing:} Attached pneumatic tubing and gradually increased internal pressure in 5 psi increments.
\end{enumerate}

\section{Observations}
% Use this section for timestamped notes as the experiment progresses.
\begin{itemize}
    \item \textbf{14:00:} Vacuum pump achieved steady state. Noticeable reduction in trapped air bubbles.
    \item \textbf{18:15:} Demolding process complete. Slight flashing observed along the parting line; trimmed with a scalpel.
    \item \textbf{18:30:} Actuator exhibits a slight curve at rest before pressure is even applied.
\end{itemize}

\section{Data \& Results}
% Table for recording raw data
\begin{table}[h]
    \centering
    \caption{Pressure vs. Bending Angle}
    \vspace{0.2cm}
    \begin{tabular}{@{} c c c @{}}
        \toprule
        \textbf{Trial} & \textbf{Pressure (psi)} & \textbf{Bending Angle ($^\circ$)} \\
        \midrule
        1 & 0 & 5 \\
        2 & 5 & 22 \\
        3 & 10 & 45 \\
        4 & 15 & 78 \\
        \bottomrule
    \end{tabular}
    \label{tab:data}
\end{table}

\section{Troubleshooting \& Next Steps}
\begin{tcolorbox}[colback=red!5!white,colframe=red!75!black,title=Issue Encountered]
At 20 psi, a micro-leak developed at the tubing insertion point, preventing further actuation.
\end{tcolorbox}
\textbf{Resolution Plan:} Next iteration will require integrating a thicker silicone collar around the inlet or using a barbed fitting cast directly into the elastomer.

\end{document}